\chapter*{译者说明}
\addcontentsline{toc}{chapter}{译者说明}
\markboth{译者说明}{译者说明}

\textbf{底本。}本书译自 Dennis Gaitsgory 与 Nick Rozenblyum 的 \emph{Crystals and D-modules}（arXiv:1111.2087v4 [math.AG]，2014 年 10 月 2 日，共 65 页）。这是一个尚在展开的研究纲领的第一部分：作者在导出代数几何的框架里定义晶体，并把它与 D-模这一经典对象联系起来。第二部分（文中记作 [GR2]）当时尚未成文。

\textbf{编排。}译文保留原书的全部编号：章的编号、节的编号（如 2.3），以及节内逐条编号的段落与定理（如 2.3.1、2.3.2）。原书中编号段与定理共用一套编号，本书照此办理。公式编号（如 (0.1)、(2.5)）亦照原书。交叉引用一律指向这些原书编号，读者可与英文原文逐条对照。原书目录中的页码是 arXiv 版的页码，本书重排后不再保留；正文中「见 §2.3」一类指向的是节号，不受影响。

\textbf{术语。}术语的选定以原书用法为准，并尽量与本译丛已有的 D-模译法保持一致：holonomic 作「完整」，direct image 作「直像」，inverse image 作「逆像」，characteristic variety 作「特征簇」，good filtration 作「好滤过」。几个关键译名在此说明：crystal 作「晶体」，prestack 作「预栈」，gerbe 与 torsor 保留英文（首现处加注），twisting 作「扭化」，eventually coconnective 作「最终余连接的」，almost of finite type 作「几乎有限型的」。全书术语见书末附表。

\textbf{译注。}正文中以灰色小字排出的段落是译者所加，共三类：易错处的预警、原文含混处的点拨、背景与直觉的补充。译注不代替原文，也不改写原文的论证；凡涉及数学内容的补充，均注明依据。来源分三级：A 级为原书、作者本人论文、经典教材与原始论文；B 级为 nLab、Stacks Project、ProofWiki 等白名单资源（须与 A 级交叉印证）；I 级为本译丛既有的译法。凡无外部依据、仅由上下文推断者，一律标明「译者按（上下文推断，无外部依据）」。原文的排印错误（如拼写错误）径行改正，不另出注；凡可能影响数学理解的笔误，则保留原样并以译注说明。

\textbf{公式。}译文中的公式逐符号对照原书页面图像核校，而不依赖文本提取。原书的交换图以 Ti\emph{k}Z 重绘，位置与形状尽量与原书一致。

\vspace{1em}
\noindent 数译\\
2026 年 11 月
\clearpage
